% Sources: original/9.4, photographs 9--16.
\chapter{Countable Sets and Number Series}

\section{Countable Sets}

\begin{notationenv}
\[
    \mathbb N=\{0,1,2,\ldots\},
    \qquad
    \mathbb N_+=\{1,2,3,\ldots\}.
\]
We also use
\[
    \mathbb Z,\qquad \mathbb Q,\qquad \mathbb R,
    \qquad
    \mathbb R_{\ge0},\qquad \mathbb R_{>0}.
\]
A sequence in a set $S$ is a map
\[
    s:\mathbb N\to S,\qquad n\mapsto s_n,
\]
denoted by $(s_n)_{n\in\mathbb N}$ or simply $(s_n)$.
\end{notationenv}

\begin{definitionenv}
A set $S$ is
\begin{enumerate}[label=(\arabic*)]
    \item \textbf{finite} if there exist $n\in\mathbb N$ and a bijection
    \[
        \{0,\ldots,n-1\}\longrightarrow S;
    \]
    \item \textbf{infinite countable} if there exists a bijection
    \[
        \mathbb N\longrightarrow S;
    \]
    \item \textbf{countable} if it is finite or infinite countable;
    \item \textbf{uncountable} if it is not countable.
\end{enumerate}
\end{definitionenv}

\begin{definitionenv}
Two sets $S_1,S_2$ are \textbf{equinumerous} if there exists a bijection
$S_1\to S_2$.
\end{definitionenv}

\begin{theoremenv}
$\mathbb N\times\mathbb N$ is infinite countable.
\end{theoremenv}

\begin{proofenv}
For $m\in\mathbb N$, let
\[
    S_m:=\{(p,q)\in\mathbb N\times\mathbb N:p+q=m\}.
\]
Enumerating $S_0,S_1,S_2,\ldots$ successively gives an enumeration of
$\mathbb N\times\mathbb N$.
\end{proofenv}

\begin{exerciseenv}
Fill in the details of the diagonal enumeration above.
\end{exerciseenv}

\begin{theoremenv}
Every subset of a countable set is countable.
\end{theoremenv}

\begin{proofenv}
It suffices to consider $A\subseteq\mathbb N$. If $A$ is infinite, define
\[
    A_0=A,\qquad a_0=\min A_0,
\]
and recursively
\[
    A_{n+1}=A_n\setminus\{a_n\},
    \qquad
    a_{n+1}=\min A_{n+1}.
\]
Then
\[
    \mathbb N\to A,\qquad n\mapsto a_n,
\]
is a bijection.
\end{proofenv}

\begin{theoremenv}
Let $S\neq\varnothing$. The following are equivalent:
\begin{enumerate}[label=(\Alph*)]
    \item $S$ is countable;
    \item there exists a surjection $\mathbb N\to S$;
    \item there exists an injection $S\to\mathbb N$.
\end{enumerate}
Moreover, $\mathbb N$ may be replaced by any infinite countable set.
\end{theoremenv}

\begin{proofenv}
$(A)\Rightarrow(B)$ is immediate for infinite $S$; for finite $S$, repeat
values of a finite enumeration.

$(B)\Rightarrow(C)$: if $f:\mathbb N\to S$ is surjective, define
\[
    g(x):=\min f^{-1}(\{x\}).
\]
Then $g:S\to\mathbb N$ is injective.

$(C)\Rightarrow(A)$ follows since every subset of $\mathbb N$ is countable.
\end{proofenv}

\begin{theoremenv}
Let $m\ge2$ and let $S_1,\ldots,S_m$ be countable
(respectively, infinite countable). Then
\[
    \prod_{j=1}^{m}S_j
\]
is countable (respectively, infinite countable).
\end{theoremenv}

\begin{proofenv}
Reduce to $m=2$ and use the countability of $\mathbb N\times\mathbb N$.
The general case follows by induction.
\end{proofenv}

\begin{exampleenv}
The set $\mathbb Q$ is infinite countable: $\mathbb N\subseteq\mathbb Q$, and
\[
    \mathbb Z\times\mathbb N_+\to\mathbb Q,
    \qquad
    (p,q)\mapsto \frac pq,
\]
is surjective.
\end{exampleenv}

\begin{theoremenv}
Let $J$ be countable and let $(S_j)_{j\in J}$ be a family of countable sets.
Then
\[
    \bigcup_{j\in J}S_j
\]
is countable.
\end{theoremenv}

\begin{exerciseenv}
Prove the preceding theorem.
\end{exerciseenv}

\begin{definitionenv}
Let $S$ be a set and $A\subseteq S$. The \textbf{indicator function} of $A$ is
\[
    \mathbf 1_A:S\to\{0,1\},
    \qquad
    \mathbf 1_A(x)=
    \begin{cases}
        1,&x\in A,\\
        0,&x\notin A.
    \end{cases}
\]
\end{definitionenv}

\begin{propositionenv}
The map
\[
    \mathcal P(S)\longrightarrow\{f:S\to\{0,1\}\},
    \qquad
    A\longmapsto\mathbf 1_A,
\]
is a bijection.
\end{propositionenv}

\begin{theoremenv}[Cantor]
$\mathcal P(\mathbb N)$ is uncountable.
\end{theoremenv}

\begin{proofenv}
Suppose there were a surjection
\[
    \mathbb N\to\{f:\mathbb N\to\{0,1\}\},
    \qquad n\mapsto f_n.
\]
Define
\[
    f(n)=
    \begin{cases}
        0,&f_n(n)=1,\\
        1,&f_n(n)=0.
    \end{cases}
\]
If $f=f_m$, then $f(m)\neq f_m(m)$, a contradiction.
\end{proofenv}

\begin{propositionenv}
\[
    \mathbb N^{\mathbb N}
    =
    \mathbb N\times\mathbb N\times\cdots
\]
is uncountable.
\end{propositionenv}


\section{Number Series}

\begin{definitionenv}
A sequence $(a_n)$ in $\mathbb R$ is \textbf{Cauchy} if
\[
    (\forall\varepsilon>0)(\exists m\in\mathbb N)
    (\forall p,q\in\mathbb N):
    \quad
    \min\{p,q\}\ge m
    \Longrightarrow
    |a_p-a_q|<\varepsilon.
\]
\end{definitionenv}

\begin{theoremenv}
A sequence in $\mathbb R$ is Cauchy if and only if it is convergent.
\end{theoremenv}

\begin{definitionenv}
For a real sequence $(a_n)$, we write $a_n\to+\infty$ if
\[
    (\forall M\in\mathbb R)(\exists m\in\mathbb N)(\forall n\in\mathbb N):
    \quad
    n\ge m\Longrightarrow a_n>M.
\]
\end{definitionenv}

\begin{exampleenv}
\[
    \lim_{n\to\infty}n^2=+\infty.
\]
\end{exampleenv}

\begin{definitionenv}
For a sequence $(a_n)_{n\in\mathbb N}$, let
\[
    s_n:=\sum_{m=0}^{n}a_m.
\]
The series $\sum_{n\in\mathbb N}a_n$ is \textbf{convergent} if $(s_n)$
converges, in which case
\[
    \sum_{n=0}^{\infty}a_n
    :=
    \lim_{n\to\infty}\sum_{m=0}^{n}a_m.
\]
\end{definitionenv}

\begin{definitionenv}
A series $\sum_{n\in\mathbb N}a_n$ is
\begin{enumerate}[label=(\arabic*)]
    \item \textbf{absolutely convergent} if
    \[
        \sum_{n\in\mathbb N}|a_n|
    \]
    converges;
    \item \textbf{commutatively convergent} if, for every permutation
    $\phi:\mathbb N\to\mathbb N$, the rearranged series
    \[
        \sum_{n\in\mathbb N}a_{\phi(n)}
    \]
    converges and its sum is independent of $\phi$.
\end{enumerate}
\end{definitionenv}

\begin{exampleenv}
\[
    \sum_{n=1}^{\infty}\frac1n=+\infty,
    \qquad
    \sum_{n=1}^{\infty}(-1)^n
    \quad\text{diverges}.
\]
\end{exampleenv}

\begin{propositionenv}
Absolute convergence implies convergence.
\end{propositionenv}

\begin{exampleenv}
The converse is false: the alternating harmonic series
\[
    \sum_{n=1}^{\infty}\frac{(-1)^n}{n}
\]
is convergent but not absolutely convergent.
\end{exampleenv}

\begin{lemmaenv}
Let $(a_n)$ be decreasing in $\mathbb R_{\ge0}$ with $a_n\to0$.
Then $\sum_{n=0}^{\infty}(-1)^n a_n$ converges.
\end{lemmaenv}

\begin{propositionenv}
There exists a permutation
\[
    \phi:\mathbb N_+\to\mathbb N_+
\]
such that
\[
    \sum_{n=1}^{\infty}
    \frac{(-1)^{\phi(n)}}{\phi(n)}
    =+\infty.
\]
\end{propositionenv}

\begin{theoremenv}[Riemann]
Let $\sum_{n\in\mathbb N}a_n$ be a convergent series in $\mathbb R$.
The following are equivalent:
\begin{enumerate}[label=(\Alph*)]
    \item $\sum_{n\in\mathbb N}a_n$ is absolutely convergent;
    \item $\sum_{n\in\mathbb N}a_n$ is commutatively convergent.
\end{enumerate}
\end{theoremenv}
